% ==============================================================================
% Section 2: Liquidation Value (Breakdown Value) Analysis
% ==============================================================================

\clearpage
\section[Liquidation Value Analysis]{Liquidation Value (Breakdown Value) Analysis}

\subsection{Theoretical Foundation and Rationale for Goodwill Exclusion}
Under corporate valuation principles, the \textbf{Liquidation Value (Breakdown Value)} represents the estimated net cash surplus realized if an enterprise were to terminate operations immediately, liquidate all physical assets through orderly secondary transactions, and extinguish all senior liabilities.

\begin{formulabox}[Governing Equation: Liquidation Value per Share]
\begin{equation}
    \LV = \frac{\text{Total Assets} - \text{Recorded Goodwill} - \text{Total Liabilities}}{\text{Equity Shares Outstanding}}
\end{equation}
\end{formulabox}

Goodwill is strictly deducted from the asset base because it represents capitalized historical acquisition premiums or intangible supernormal earning power. Under distress or bankruptcy conditions, commercial goodwill rapidly dissipates, possessing negligible liquidation recovery value.

\subsection{Cross-Sectional Empirical Results}
Table~\ref{tab:liquidation_value} summarizes the balance-sheet breakdown metrics, net realizable values, and resulting liquidation floors for the ten sampled FMCG corporations.

\begin{table}[H]
\centering
\footnotesize
\setlength{\tabcolsep}{2.0pt}
\renewcommand{\arraystretch}{1.08}
\begin{tabular*}{\textwidth}{@{\extracolsep{\fill}} l r r r r r r r}
\toprule
\textbf{Company} & \shortstack[r]{\textbf{Total}\\\textbf{Assets}\\\scriptsize(\inr Cr)} & \shortstack[r]{\textbf{Goodwill}\\\scriptsize(\inr Cr)} & \shortstack[r]{\textbf{Total}\\\textbf{Liab.}\\\scriptsize(\inr Cr)} & \shortstack[r]{\textbf{Net Realiz.}\\\textbf{Value} (\inr Cr)} & \shortstack[r]{\textbf{Shares}\\\scriptsize(Cr)} & \shortstack[r]{\textbf{LV / Share}\\\scriptsize(\inr)} & \shortstack[r]{\textbf{Market}\\\textbf{Price} (\inr)} \\
\midrule
\textbf{Hindustan Unilever}   & 79,738 & 1,478 & 30,999 & 47,261 & 235   & \textbf{\inr 201} & \inr 1,973 \\
\textbf{ITC Ltd}              & 93,637 & 2,399 & 16,981 & 74,257 & 1,253 & \textbf{\inr 59}  & \inr 264   \\
\textbf{Nestle India}         & 13,357 & 444   & 6,641  & 6,272  & 193   & \textbf{\inr 33}  & \inr 1,409 \\
\textbf{Varun Beverages}      & 25,541 & 0     & 7,639  & 17,902 & 676   & \textbf{\inr 26}  & \inr 204   \\
\textbf{Britannia Industries} & 9,730  & 1,380 & 4,648  & 3,702  & 24    & \textbf{\inr 154} & \inr 5,120 \\
\textbf{Marico Ltd}           & 9,950  & 557   & 5,183  & 4,210  & 130   & \textbf{\inr 32}  & \inr 814   \\
\textbf{Tata Consumer}        & 34,279 & 2,820 & 11,264 & 20,195 & 99    & \textbf{\inr 204} & \inr 1,010 \\
\textbf{Godrej Consumer}      & 13,365 & 2,948 & 5,573  & 4,844  & 102   & \textbf{\inr 47}  & \inr 881   \\
\textbf{Dabur India}          & 17,480 & 1,287 & 6,238  & 9,955  & 177   & \textbf{\inr 56}  & \inr 382   \\
\textbf{Colgate-Palmolive}    & 3,408  & 47    & 1,394  & 1,967  & 27    & \textbf{\inr 73}  & \inr 1,843 \\
\bottomrule
\end{tabular*}
\caption{Cross-Sectional Liquidation Value (Breakdown Value) Valuation Summary.}
\label{tab:liquidation_value}
\end{table}

\subsection{Analytical Interpretation and Valuation Multiples}
\begin{enumerate}[label=\textbf{\arabic*.}, topsep=2pt, itemsep=1.5pt, parsep=0pt]
    \item \textbf{Universal Overvaluation Relative to Breakdown Floor:} All ten firms trade at substantial multiples above liquidation floors, ranging from \textbf{$4.5\times$} (ITC Ltd) to an extreme \textbf{$42.7\times$} (Nestle India Ltd).
    \item \textbf{Capital Structure and Downside Protection:} \textbf{ITC Ltd} ($P_0 = \inr 264$, $\LV = \inr 59$) and \textbf{Tata Consumer} ($P_0 = \inr 1,010$, $\LV = \inr 204$, $4.9\times$ multiple), backed by \inr 74,257\crunit\ and \inr 20,195\crunit\ in net tangible assets respectively, offer the strongest balance-sheet cushion against solvency shocks.
    \item \textbf{Asset-Light Franchise Premium:} Formulators like \textbf{Nestle India} ($42.7\times$) and \textbf{Britannia} ($33.2\times$) command immense multiples because tangible plant and machinery generate only a fraction of cash flows; value is driven by uncapitalized brand equity, pricing power, and distribution moats.
\end{enumerate}
